\documentclass{report}
\usepackage{amsmath,amssymb}
\setlength{\parindent}{0mm}
\setlength{\parskip}{1em}
\begin{document}
\begin{center}
\rule{6in}{1pt} \
{\large J. Zhao \\
{\bf A multigrid method with an FFT smoother and deflation techniques for elastic contact problems}}

HB03 060 \\ Mekelweg 4 \\ 2628 CD Delft \\ The Netherlands
\\
{\tt J.Zhao-1@tudelft.nl}\\
E.A.H. Vollebregt\\
C.W.  Oosterlee\end{center}

In the simulation of railway vehicles dynamics, the interaction between
the vehicles' wheels and rails attracts a lot of interest. It involves
the solution of the so-called {\it contact problems}, concerning the
normal and tangential tractions on the inter-surface. The formulation
based on Kalker's variational half-space approach is regarded as an
accurate model for contact problems, particularly those involving a
rolling contact with friction. Fast solvers are demanded for such
problems.


In this talk, multigrid methods are used to solve the discretized
governing system, which is derived from an integral equation and its
coefficient matrix $A$ is full, symmetric and Toeplitz. A smoother is
proposed, which is based on the Richardson iteration, where the residual
in each iteration is preconditioned by a matrix $M^{-1}$. The idea of
constructing this matrix comes from the fact that each column of the
inverse of the coefficient matrix $A^{-1}$ is similar, and we approximate
any one of these by a fast Fourier transform (FFT) technique. The
resultant vector defines a Toeplitz matrix which is the preconditioner
$M^{-1}$.


This smoother possesses the advantages of easy construction, and fast
computation of matrix-vector multiplications using FFT. It is able to
eliminate the high frequency modes of the errors to a great extent,
however, it also enlarges several very low frequency modes. This is
remedied by a subdomain deflation approach, where only a few
piecewise-constant deflation vectors are required.


The numerical tests and some spectral analysis indicate very fast
convergence and mesh-independence of our method. Hence this method may
accelerate the solution procedure for the contact problems. In view of
other applications, on the one hand, it can be applied for systems with
Toeplitz coefficient matrices arising in other engineering fields such as
image processing. On the other hand, a preconditioner based on these
techniques also performs well in combination with Krylov subspace methods
aiming at even faster convergence.


\end{document}
